% =====================================================================
%  sec02-established.tex
%  Part II: what is established.
%  Sources read for this section:
%    research/v3/lean/{M02-CBAR, L04PropC, L05Counterexamples, M03-COUNTER}.lean
%    research/v3/species/{C06-WINDOW, ADJ2, ADJ4, C07-AUDIT, F07-ADJCBAR,
%      D05-FKG, N09-RESTATE, L02-CBARKATO, L06-INTERVAL, Y01-KL,
%      Z06-HOLEA, Z07-HOLEB, Z08-HOLEC, Z09-HOLED, Z10-HOLEE, Z12-FIXC}.md
%    docs/paper-v2/ch20-unconditional-i1.tex (five rings, five holes, window)
%    docs/paper-v7/sg1/main.tex and docs/paper-v7/w1/main.tex
%    docs/position-current.html:373--417 (round-178 status block)
%  House rules honoured here: every number carries a source file; grades are
%  never written above what the sources support; no completion stamp is
%  printed anywhere; Clay = 0%.
% =====================================================================
\clearpage
\begingroup
\renewcommand{\code}[1]{\texttt{\footnotesize\detokenize{#1}}}

\part{What is established}


\section{What is established: the residue, the isolation, the uniqueness, the no-go, and four machine-verified lemmas}
\label{sec:established}

\noindent This section is the register of statements that survived audit. It is deliberately narrow:
every claim below carries a grade chip and the file that produced it, and no number is quoted from
memory. Of the six subsections, three are theorems in the strict sense (\LT: $I1$, $I2$, and the
block-Loewner order), two are theorems under hypotheses printed with them (\LC: $N09$-$1$ and
$SG1$-$1$), one is model-level (\LM: the four machine-verified lemmas), and two carry corrections that
reduce what an earlier edition of this work asserted --- the window of \S\ref{sec:est-window-fix} and
the grade of \S\ref{sec:est-sg1-grade}.

\smallskip
\noindent\textbf{Two grading rules, used throughout.} (1) A claim is graded by the \emph{weakest}
link of the chain that certifies it, not by the strength of its conclusion. (2) A \emph{computed}
number never upgrades a \emph{proved} structure, and a proved structure never upgrades a computed
number: in \S\ref{sec:est-i1} the two are physically separated inside one theorem, and that
separation is the point of the subsection. \src{\code{position-current.html:373--380}.}

\noindent\textbf{Scope, stated once and not repeated.} Nothing in this section crosses to
$\beta\approx2.5$, nothing takes $a\to0$, and no statement below mentions the continuum mass gap.
The residue theorem is a strong-coupling statement at fixed spacing; its window is in the coupling,
and the spacing has a positive lower bound inside it (\code{docs/paper-v2/ch20-unconditional-i1.tex}
\S D.3). \textbf{Clay remains $0\%$.}

% =====================================================================
\subsection{The strong-coupling residue theorem \texorpdfstring{$I1$}{I1}: \texorpdfstring{\LT}{[T]} on the structure, \texorpdfstring{\LC}{[C]} on the value}
\label{sec:est-i1}

\begin{claimbox}{cT}
\claimhead{Theorem}{(strong-coupling residue $I1$).} Let
$\mathcal C_L=T^3_L\times_a\mathbb Z$ be the cylinder (spatial torus of side $L$, discrete time of
spacing $a$), let $O=\tfrac14\operatorname{tr}F^2$, and let $m_1$ be the lightest energy carrying
nonzero $O$-overlap. Fix $N$ and work in the SZZ coupling $\beta_A=2/(g^2N)$ (Wilson:
$\beta_W=N^2\beta_A$). For every finite $L$, \emph{uniformly in $L$},
\begin{equation}
\beta_A\ <\ \beta_\star\ :=\ W(a^\ast)\ =\ 3.777\times10^{-3}\ =\ 0.1813\times\tfrac1{48},
\qquad\text{i.e.}\qquad
\beta_W\ <\ 0.015108\ =\ 0.725\times\tfrac1{48},
\label{eq:est-window}
\end{equation}
the lightest $0^{++}$ atom of $O$ obeys
\begin{equation}
f^2_{1,L}\ :=\ \nu_L(\{E_{1,L}\})\ =\ \lim_{t\to\infty}e^{E_{1,L}t}C_L(t)\ \ge\ r_0\ >\ 0 .
\label{eq:est-f1}
\end{equation}
The order of limits is: $L$ finite and fixed, then $t\to\infty$; $L$ is a uniformity axis and
$L\to\infty$ is never taken. \emph{Structure: \LT. The window \eqref{eq:est-window}: \LC, computed
and not proved} (\S\ref{sec:est-window-fix}).
\end{claimbox}

\noindent\textbf{The five rings, all \LT.} The chain to \eqref{eq:est-f1} has five links. Each is one
line here because the work is in the file; naming them is what makes the grade auditable.

\begin{enumerate}[label=\textbf{R\arabic*.},leftmargin=2.7em]
\item \textbf{Four walls} (\code{P3-FOURWALL-DERIVATION.md} \S1.1--1.5). Each of the source
      plaquette's four boundary links carries an open fundamental index pair and $\int dU\,U_{ij}=0$
      forbids covering it alone; two unit plaquettes share at most one link, so there are at least
      four independent strands; only temporal plaquettes advance time, so each strand crosses $t/a$
      slices. Hence $|P|\ge4t/a$, with equality exactly for the straight four-walled tube, whose
      weight is $a_0u^{4t/a}=a_0e^{-m_1t}$ with $m_1=4\ln(1/u)/a$. \LT.
\item \textbf{KP in the full representation} (\code{P3-FULLREP-KP.md} \S2.1--2.3, \S4.1). The
      criterion is Koteck\'y--Preiss \cite{kp}. Haar is a
      probability measure, so $|\chi_r(U)|\le\dim r$; the multiplicity cap
      $c_r\le\beta^{|r|}/(|r|!\dim r)$ cancels exactly the $\dim r$ that the contraction would cost;
      the factorial beats the sub-exponential rep-label entropy; hence $\sum_{P\ni x}|w(P)|e^{a|P|}
      \le a$ whenever $uK<1$, with $K\le6.1\times10^2N$ and plaquette adjacency
      $\Delta=4(6-1)=20$. \LT.
\item \textbf{Resummation of the decorations} (\code{X01-L1C.md} \S1--2, \code{Y04-AUDITX01.md}
      \S1--3; activities corrected in \code{Z08-HOLEC.md}, \code{Z12-FIXC.md}). Clusters are
      grouped by attachment span $\alpha$, not by support; the tube is open, so placements are
      counted by $(T-\alpha+1)_+$; the shift $\alpha\mapsto\alpha-1-T$ moves weight and exponent
      together, so the tail is $|\varepsilon_T|\le B\bar u^T$ with no $T$-prefactor; the
      interface-corrected activity has envelope $\hat C=4$. \LT.
\item \textbf{Spectral bridge} (\code{P3-CYLINDER-RESIDUE.md} \S1.2--1.4). At finite $L$ the
      transfer matrix is Hilbert--Schmidt on a compact configuration space, hence compact, hence
      purely discrete: every spectral point is an atom, and $f_1^2=\lim_te^{E_{1,L}t}C_L(t)$ uses
      only $\delta_L=E_{2,L}-E_{1,L}>0$ at that $L$ \cite{reed-simon}. \LT.
\item \textbf{$[KL]$} (\code{P3-CYLINDER-RESIDUE.md}:67--69, \code{Y01-KL.md} \S4.1).
      $\nu_L\ge0$ is Osterwalder--Seiler reflection positivity \cite{osterwalder-seiler,os1973} and
      $f_{n,L}^2\ge0$ is the spectral theorem, so the K\"all\'en--Lehmann representation \cite{kl} on
      the cylinder is a theorem with no Yang--Mills dynamics in it, and it transfers verbatim to the
      $2^{++}$ channel. \LT.
\end{enumerate}

\noindent\textbf{The five holes A--E, all closed.} Each hole is a specific defect that an earlier
draft absorbed by relabelling; each is now closed in its own file.

\begin{table}[htbp]
\centering
\footnotesize
\setlength{\tabcolsep}{4pt}
\renewcommand{\arraystretch}{1.2}
\begin{tabularx}{\textwidth}{@{}l>{\raggedright\arraybackslash}p{3.1cm}%
  >{\raggedright\arraybackslash}X>{\raggedright\arraybackslash}p{3.3cm}@{}}
\toprule
\textbf{Hole} & \textbf{What it was} & \textbf{How it is closed} & \textbf{Source} \\
\midrule
A & $u=c_f$ vs.\ $u=c_f/N$: a factor $N^{4T}$ in the tube weight
  & $u:=c_f/N$, one $1/N$ per tile, from Gupta (14.18)/(14.21); and $A^{\mathrm{dressed}}$
    \emph{is} $f_1^2$, so no second quantity ever existed
  & \code{Z06-HOLEA.md} \S1--5 \LT \\
B & whether $\delta\ge\kappa/3$ lands inside $I1$
  & out: neither $I1$ nor $D_\infty$ contains $\delta$; the reachable rate is
    $\delta=\ln(1/\bar u)/a=\kappa/4-\ln(\hat Ce\Delta\bar C)/a$, strictly below $\kappa/4$
  & \code{Z07-HOLEB.md} \S2--4 \LT \\
C & re-coupling is not block-multiplicative: $\mathrm{Hom}(1,f\bar ff\bar f)=2\ne1^2$; caps have
    $w(\gamma)=0$ while $w(B_T\cup\gamma)\ne0$
  & $\hat C=2$ exactly at each factorization point, envelope $\hat C=4$ per plaquette; the
    interface re-coupling is absorbed into the activity, never into $\varepsilon_T$
  & \code{Z08-HOLEC.md} \S2--4, \code{Z12-FIXC.md} \S1--3 \LT \\
D & summing $O$ over the slice makes $f_1^2\propto L^3$, so a bare $f_1^2$ cannot meet an
    $L$-independent $r_0$
  & $r_0$ carries the volume factor (equivalently $f_1^2$ is read normalized); $a_0$'s
    $L$-independence is already proved by a positive geometric sum with wrap corrections $e^{-cL}$
  & \code{Z09-HOLED.md} \S3--5 \LT \\
E & $r_0=\tfrac14(2N)^{-2}$ vs.\ $\tfrac12c_1u^p$: two $\beta$-scalings for one constant
  & $p$ is the operator-coefficient exponent and $p=0$; $u$ lives only in the correction factor,
    $r_0=\tfrac12a_0(1-O(u))$, $\beta$-independent
  & \code{Z10-HOLEE.md} \S3--5 \LT \\
\bottomrule
\end{tabularx}
\caption{The five holes A--E of the residue-and-isolation chain, each with the defect that it was and
the argument that closes it. Every row is \LT{} in its own file; the sources are named in the last
column.}
\label{tab:holes}
\end{table}

\noindent\textbf{The two accounting constants.} \emph{$\bar C\equiv N=2$, as an identity and not as
an upper bound.} The Ursell constant is $\bar C=c_f/u=N$, one $1/N$ per link from the integration,
with $|w|^{1/n}/u=c_f/u=N$ on cubes and four-walled tubes (measured $2.0000$ at $n=6,4,8,12$); the
$\mathrm{SU}(2)$ cross-check at $\beta=1/12$ is $\langle\tfrac1N\operatorname{tr}U_p\rangle
=0.0208273$ against $c_f=0.0208333$, a $0.03\%$ gap.
\src{\code{F07-ADJCBAR.md} \S3; \code{M02-CBAR.lean}, \code{cbar_eq_N}, \code{cbar_eq_N_two}.}

\smallskip
\noindent\emph{$\hat C=4$, and attained rather than bounded.} The KP criterion sums over polymers
$\gamma$, so what enters is the per-polymer rate $\sup_\gamma(\prod_\ell R(\ell))^{1/|\gamma|}$; the
handshake identity $E_2/F=2$ on closed clusters gives $2^{E_2/F}=2^2=4$, and a witness attains it:
all $324/324$ cubes of $T^4_3$ give geometric mean $4.000000$, as do $4000$ random closed
complexes, with $0$ violations. The "$2$" is the value at each factorization point, not the value on
a polymer, so it can never replace the $4$. \LT{}$+$\LE.
\src{\code{ADJ2.md}:13; \code{C06-WINDOW.md} \S2.2--2.3.}

\smallskip
\noindent\emph{The two constants meet the threshold.} $\hat C\bar C=8$ against
$2N^2\cdot48/e\Delta=7.0632853\ldots$, and the Lean file contains the integer certificate for
$8>7.0632853$. \src{\code{M02-CBAR.lean}: \code{chat_times_cbar}, \code{thresh},
\code{thresh_iff}, \code{thresh_eq_eight_div_x}.}

\smallskip
\noindent\emph{Penrose, approved with its own caveat.} The tree bound $|\varphi^U(\Gamma)|\le
\tau(G_\Gamma)$ is \LT{} with a four-line deletion--contraction proof plus $875$ graphs and $0$
violations; the \emph{citation} to Penrose (1967) and to the Penrose identity in
Fern\'andez--Procacci stays \LC{} because the original page was never opened
(\code{C07-AUDIT.md}:17, :24, :124--126, :135).

\smallskip
\noindent\emph{Lean.} The arithmetic core of the two constants is $25$ theorems in
\code{M02-CBAR.lean} with \code{sorry} $=0$ and twelve axiom-audit lines returning
$[\code{propext},\code{Classical.choice},\code{Quot.sound}]$; the last run exits $0$ with no error
and no warning. \src{\code{M02-CBAR.lean} (theorem count and axiom audit at its end);
\code{LOOP-RUNBOOK.md}:9512--9528.}

% ---------------------------------------------------------------------
\subsubsection{The window correction: the value drops to \texorpdfstring{\LC}{[C]}, the structure does not drop}
\label{sec:est-window-fix}

\begin{warnbox}
\textbf{What was over-reported, by how much, and why the old number is illegal.} The window printed
in the earlier edition of this chapter was $\beta<0.018394$. That number is
$1/(e\Delta)$, the \emph{activation root at $a=0$} of the KP criterion --- and $a=0$ sets the
right-hand side of the criterion to zero, so the criterion is identically false there and the root
is a convergence radius, not a criterion window. \textbf{It is therefore not a legal window.} With
the grading parameter optimized,
\begin{equation}
\beta_A\ <\ W(a^\ast)\ =\ 3.777\times10^{-3}\ =\ 0.1813\times\tfrac1{48},
\qquad\Longleftrightarrow\qquad
\beta_W\ <\ 0.015108\ =\ 0.725\times\tfrac1{48},
\label{eq:est-window2}
\end{equation}
so the earlier report over-states the window by
\[
\frac{0.018269}{0.003777}\ =\ 4.84\times ,
\qquad\text{and this chapter's own}\quad
\frac{0.018394}{0.003777}\ =\ 4.87\times .
\]
\src{\code{ADJ4.md}:16 (the ratio $1:2.7508:7.7238$ and $\Phi$ table), :19 and :125 (the two
over-report factors); \code{LOOP-RUNBOOK.md}:17133--17194 (\S131, the same two factors);
\code{C06-WINDOW.md} \S2.}
\end{warnbox}

\noindent\textbf{The optimizer, printed so that it can be checked.} With $\lambda=e\Delta=54.3656$
and $\Phi$ the single-plaquette sum, the criterion $S_x(a)\le a$ is equivalent to
$\lambda\Phi(a,\beta)\le a/(1+a)$, hence $W(a)=a/[(1+a)\lambda\Phi(a)]$. The measured ratios are
$\Phi(0):\Phi(1):\Phi(2)=1:2.7508:7.7238$ at $\beta=0.018394$ ($2.75512$ at $\beta=1/48$), and the
numerical maximum sits at $a^\ast=0.6150$. The closed-form point $a^\ast=(\sqrt5-1)/2=0.618$ is an
approximation to that maximum, with the same value to four digits, and $a=1$ gives
$3.372\times10^{-3}$ (the $3.357\times10^{-3}$ of the earlier audit, $0.44\%$ lower).
\src{\code{ADJ4.md}:16, :19; \code{CB4.md}:53--57 as quoted there; \code{LOOP-RUNBOOK.md}:17149.}

\smallskip
\noindent\textbf{And the natural coupling is outside, by a wide margin.} At
$\beta_W=1/12=0.08333$ one has $u=\beta_W/(2N^2)=0.0104$ and $\bar u=e\Delta N u=1.132617>1$, so the
criterion fails; against the corrected window the same coupling is out by
$0.08333/0.015108=5.5\times$. \src{\code{M02-CBAR.lean} (\code{cf} at $\beta=1/12$ gives
$c_f=1/48$); \code{BETA-CONVENTION.md}:187; \code{F07-ADJCBAR.md} \S4.}

\begin{qbox}
\textbf{Why this is a correction of a value and not a retraction of a theorem.} The window is one
slot of the statement \eqref{eq:est-window} and it is not a premise of any ring of the chain: the
five rings \LT{} deliver $f^2_{1,L}\ge r_0>0$ for every $\beta$ at which the KP criterion holds, and
``at which it holds'' is exactly what the window names. Re-optimizing the grading parameter changes
\emph{where} the criterion holds, not \emph{what} it delivers. Hence the grade splits: the value of
the window is \LC{} (computed: a numerical maximum of an explicit function), the residue statement
stays \LT. \textbf{There is no conflict between the correction and the \LT{} structure, because the
two occupy different slots of the same sentence}; writing the whole theorem \LC{} would under-state
the chain, and writing the window \LT{} would over-state a numerical maximum.
\src{\code{ADJ4.md}:19 (the verdict separates the number from the reason); \code{position-current.html}:375.}
\end{qbox}

% =====================================================================
\subsection{The isolation \texorpdfstring{$I2$}{I2} and Proposition C}
\label{sec:est-i2}

\begin{claimbox}{cT}
\claimhead{Theorem}{(isolation $I2$).} In the setting of \S\ref{sec:est-i1}, with
$\Lambda_\star:=\inf\operatorname{supp}\,(\nu-f_1^2\delta_{m_1^2})$ the bottom of the spectrum above
the leading atom, for every finite $L$ and uniformly in $L$,
\begin{equation}
\Lambda_\star-m_1^2\ \ge\ \delta\ >\ 0 ,
\label{eq:est-i2}
\end{equation}
and the same holds in the $2^{++}$ channel with its own $\nu^{(2)}$, $m_2$ and
$\Lambda^{(2)}_\star$. \LT.
\end{claimbox}

\noindent\textbf{Why the uniformity in $L$ is the hard half.} Isolation at one $L$ is a statement
about one transfer matrix; isolation uniformly in $L$ is a statement about a one-parameter family,
and gaps along a family are exactly what degeneracies destroy: the von Neumann--Wigner phenomenon
\cite{vnw} makes a crossing generic once the symmetry that protects it is broken. The five rings earn
the uniformity; an estimate at fixed $L$ would not.

\noindent\textbf{The five rings are the same five.} Isolation is not a second chain: rings R1--R5 of
\S\ref{sec:est-i1} certify \eqref{eq:est-f1} and \eqref{eq:est-i2} on the same transfer matrix, and
the five holes A--E are the same five holes. This is why the two statements share a subsection in
the corpus and are separated here only for readability.
\src{\code{docs/paper-v2/ch20-unconditional-i1.tex}, its five-ring section and its five-hole table;
\code{Y02-CHANNELS.md}:16--20.}

\smallskip
\noindent\textbf{Proposition C, machine-verified.} The isolation has a finite-dimensional shadow that
is a Lean theorem. In the finite-support model $\mathrm{Spec}=(\text{points},\text{weights})$, put
$I2:=\exists\Lambda\,(\Lambda$ is the least support point above $m_1$ and $m_1^2<\Lambda)$ and
$\mathrm{SecExists}:=\exists x\in\text{points},\,m_1<x$. Then
\begin{equation}
0<m_1<1\quad\Longrightarrow\quad\bigl(I2\ \Longleftrightarrow\ \mathrm{SecExists}\bigr),
\label{eq:est-propc}
\end{equation}
with \code{sorry} $=0$. The forward direction is unconditional; the backward direction is exactly
where the window $0<m_1<1$ is spent, through $m_1^2<m_1$.
\src{\code{L04PropC.lean}: \code{propC}, \code{propC_fwd}, \code{propC_bwd}, \code{sq_lt_self};
\code{LOOP-RUNBOOK.md}:9109--9129.}

\smallskip
\noindent\textbf{Two limits of Proposition C, recorded because they are load-bearing.}
\emph{(i)} $I1$ is automatic only when the point table is declared to be the positive-weight support;
with a bare grid of points $I2$ does not imply $I1$, and the file contains that counterexample
(\code{counterexample_I2_not_I1}). \emph{(ii)} The convexity lemma $A1$ is machine-verified in its
finite-sum form only: \code{MeasureTheory} is unavailable in the build, so the integral form stays
\LC{} (\code{L04PropC.lean}, header and \code{a1}; \code{LOOP-RUNBOOK.md}:9129). Neither limit
touches \eqref{eq:est-i2}, which is not a finite-support statement.

% =====================================================================
\subsection{Uniqueness of the target: \texorpdfstring{$N09$-$1$}{N09-1}}
\label{sec:est-n09}

\begin{claimbox}{cC}
\claimhead{Theorem}{(uniqueness of the target, $N09$-$1$; conditional on the M6 specification and
the $L\to\infty$ reading).} Let a quantity $X$ be required to pass the four-model slot test --- free
field $0$, $2$D Yang--Mills ``true'', M6 ``false'', single plaquette ``undefined'' --- with the four
slots kept distinct. Then $X$ is a function of the mass: there is $f$ with $f(0)=0$ and $f>0$ for
$m>0$ such that $X=f(m)$. \LC. Consequently the target is unique: every admissible ``new
invariant'' is a relabelling of $m$.
\end{claimbox}

\noindent\textbf{The grade.} The two hypotheses --- the M6 specification and the $L\to\infty$ reading
of its spectrum --- are printed in the statement, and by the grading rule of the title page a theorem
carrying printed hypotheses is \LC, not \LT. The earlier edition graded this statement \LT; that is
the same species of over-report as the window corrected in \S\ref{sec:est-window-fix}, and it is
corrected here rather than quietly. The content is unaffected: the hypotheses are the ones under
which $M6$ was constructed, and no use of $N09$-$1$ in this paper leaves them.

\noindent\textbf{Why.} M6 is oscillator $\otimes$ free mode with $w=1$, and by the $L\to\infty$
reading its only vanishing scale is the spectral bottom $m_6=\min(w,2\pi^2/L^2)$. A quantity that is
defined on M6 and not identically true must take the value $0$ there, i.e.\ must vanish with
$m_6\to0$; on $2$D Yang--Mills it is positive where $m=3g^2/8>0$. A functional that is defined,
non-vacuous and zero exactly on the degenerate model is therefore a monotone relabelling of $m$ at
$0$. \src{\code{N09-RESTATE.md}:44--53 (\S2, the four numbered steps), :24--30 (four-valued
signature).}

\smallskip
\noindent\textbf{The qualification, and it is the thirty-eighth ledger entry.} Statement
$N09$-$1$ holds \emph{only at fixed operator $O$}. The four-model test fixes the \emph{model}; it
does not fix the \emph{operator}, and the atom weight is a property of the pair. On one ensemble, one
lattice, one $\beta$, with no change of $H$, the same Markov chain and the same $12000$
configurations give
\begin{equation}
\begin{gathered}
w_1=0.3427\pm0.0119\quad(O=\tfrac14\operatorname{tr}F^2,\ \text{unsmeared}),\\
w_1=0.0168\pm0.0142\quad(\text{spatial APE }n=1),
\end{gathered}
\label{eq:est-w1op}
\end{equation}
a ratio $0.343/0.017\approx20.2$. The correct reading of $N09$-$1$ is therefore
$X=f(m;O)$: uniqueness is uniqueness \emph{per operator}, and any statement of the form
``$w_1>0$'' has no content until $O$ is named.
\begin{sloppypar}
\noindent\src{\code{docs/paper-v7/w1/main.tex}:233, :251, :261--264 (the pair and the ratio
$20.2$); \code{position-current.html}:377 (entry 38); \code{ATOM2.md}:15--16 for the one-line
counterexample $w_1=\cos^2\theta$ at fixed spectrum.}
\end{sloppypar}

% =====================================================================
\subsection{The no-go \texorpdfstring{$SG1$-$1$}{SG1-1}: \texorpdfstring{\LC}{[C]}, not \texorpdfstring{\LT}{[T]}}
\label{sec:est-sg1}

\begin{claimbox}{cC}
\claimhead{Theorem}{(no-go $SG1$-$1$; grade \LC).} On the class $\mathcal M$ of non-negative Borel
measures on $[0,\infty)$ with finite moments up to a fixed $N$ and with
$s_1(\nu)=\inf\operatorname{supp}\nu\in\mathbb R$, there is no functional
$F:\mathcal M\to\mathbb R$ satisfying simultaneously
\begin{equation}
\begin{gathered}
\textbf{(i)}\ \text{$G1$-covariant},\qquad
\textbf{(ii)}\ F\le s_1,\qquad
\textbf{(iii)}\ \text{finite-moment computable},\\
\text{hence}\qquad
\textbf{(i)}\wedge\textbf{(ii)}\wedge\textbf{(iii)}\ =\ \varnothing .
\end{gathered}
\label{eq:est-sg1}
\end{equation}
\textbf{The grade is \LC{} and not \LT}, and \S\ref{sec:est-sg1-grade} states exactly which step
costs the downgrade.
\end{claimbox}

\begin{proofx}
Fix $a\ge0$, $\lambda>0$ and $0<\varepsilon<\min(a,2\lambda)$, and put
\begin{equation}
\begin{gathered}
\nu_\varepsilon^{(a)}\ =\ \frac{\varepsilon}{2\lambda}\,\delta_\varepsilon
\ +\ \Bigl(1-\frac{\varepsilon}{2\lambda}\Bigr)\,\delta_a\ \in\ \mathcal M ,
\qquad
s_1\bigl(\nu_\varepsilon^{(a)}\bigr)=\varepsilon ,\\
m_k\bigl(\nu_\varepsilon^{(a)}\bigr)\ =\ \frac{\varepsilon}{2\lambda}\,\varepsilon^{k}
+\Bigl(1-\frac{\varepsilon}{2\lambda}\Bigr)a^{k}
\ \xrightarrow[\varepsilon\downarrow0]{}\ a^k=m_k(\delta_a)
\qquad(k\ge1).
\end{gathered}
\label{eq:est-atom}
\end{equation}
\emph{Step 1.} By \eqref{eq:est-atom} the moment vector of $\nu_\varepsilon^{(a)}$ converges to that
of $\delta_a$. \emph{Step 2.} Continuity of $\Phi$ in (iii) gives
$F(\nu_\varepsilon^{(a)})\to F(\delta_a)$, while (ii) gives
$F(\nu_\varepsilon^{(a)})\le s_1(\nu_\varepsilon^{(a)})=\varepsilon$; hence $F(\delta_a)\le0$ for
every $a\ge0$. \emph{Step 3.} $\tau_a\delta_0=\delta_a$, so (i) gives
$F(\delta_a)=F(\delta_0)+a$; combined with Step 2 this forces $F(\delta_0)+a\le0$ for every $a\ge0$,
which fails for $a>|F(\delta_0)|$. No such $F$ exists. The argument uses one moment ($N\ge1$) and
continuity, and nothing else.
\end{proofx}

\noindent The three-line skeleton is the corpus statement, transcribed word for word: the atom family
drifts the moments by $O(\varepsilon)$ while $s_1:1\to\varepsilon$; the machine check at
$a=0.5/1/2/5$ returns \textsf{NO} in all four rows; hence the only escape is the full-moment limit,
the only reachable object there is $F=s_1$, so $G\equiv0$, which collides with $N09$-$1$.
\src{\code{SG1.md}:62--67 as quoted in \code{docs/paper-v7/sg1/main.tex}:252--268;
\code{w7/sg1_G.log}:37--42 (the four \textsf{NO} rows), reproduced at
\code{docs/paper-v7/sg1/main.tex}:401--404 (values $1.0\cdot10^{-6}$, $2.49999\cdot10^{-7}$,
$4.99999\cdot10^{-7}$, $9.99999\cdot10^{-7}$, $2.5\cdot10^{-6}$ against $s_1$ equal to
$0.5/1/2/5$).}

% ---------------------------------------------------------------------
\subsubsection{Why the grade of \texorpdfstring{$SG1$-$1$}{SG1-1} is \texorpdfstring{\LC}{[C]}: the thirty-seventh ledger entry}
\label{sec:est-sg1-grade}

\noindent Condition (iii) admits exactly one reading in the proof above: \emph{there is an
$N<\infty$, fixed independently of $\nu$, and a continuous $\Phi:\mathbb R^N\to\mathbb R$ with
$F(\nu)=\Phi(m_1(\nu),\dots,m_N(\nu))$ for every $\nu\in\mathcal M$}. Under that reading Step 2 is a
one-line consequence of continuity, and \eqref{eq:est-sg1} is a theorem.

\smallskip
\noindent\textbf{If (iii) is read as computability} --- recursion-theoretic, or ``computable from
$(S_W,\beta,L,a)$ by an algorithm'' --- then the passage
$F(\nu_\varepsilon^{(a)})\to F(\delta_a)$ does not follow: a functional computable in the moment
data but discontinuous in the moment topology is outside the killed class. \textbf{That single step
is therefore \LC, and with it the true grade of $SG1$-$1$ is \LC, not \LT.} The proof is three lines
long only because continuity is doing all of the work; remove continuity and nothing is left.
\src{\code{docs/paper-v7/sg1/main.tex}:207--232 (convention \code{iii-strict} and the remark on
(iii)), :320--330 (Step 2 carries the chip); ledger entry 37,
\code{position-current.html}:378.}

\smallskip
\noindent\textbf{Two conventions that the statement needs, and neither is cosmetic.}
\emph{(1) (i) is pinned to the G1 energy axis.} Under a spatial translate the object
$\inf\operatorname{supp}\nu$ is invariant, so a G2 reading makes (i) constrain nothing and
simultaneously puts it in direct conflict with (ii); the conflict is an artefact of the reading.
\emph{(2) (ii) quantifies over the whole measure class} $\mathcal M$, not merely over the spectral
measures physics supplies. That makes the emptiness \emph{stronger} than its physical-class version:
a functional that fails on the larger class fails a fortiori on the subclass.
\src{\code{B05-DUAL.md}:76--83, :80; \code{docs/paper-v7/sg1/main.tex}:163--205 (conventions
\code{domain} and \code{G1}).}

\smallskip
\noindent\textbf{Relation to $N09$-$1$: no implication in either direction.} $N09$-$1$ quantifies
over \emph{quantities on models} that pass the slot test; $SG1$-$1$ quantifies over
\emph{functionals on measures}. They meet at exactly one point, and that point is a degenerate end:
the survivor $F=s_1$ gives $G\equiv0$, which passes $4/4$ and lands inside $\operatorname{dom}(N09$-$1)$
as the function $f\equiv0$. The corpus phrase ``collides with $N09$-$1$'' therefore means: the
non-trivial defect $G$ is excluded, and the zero defect is not an object at all. It does not mean
that one statement entails the other; the two quantify over different carriers, and each is
conditional on a hypothesis the other does not share --- the $M6$ specification for $N09$-$1$, the
fixed-$N$ continuity reading of (iii) for $SG1$-$1$ --- so no grade transfers between them.
\src{\code{docs/paper-v7/sg1/main.tex}:617--641 (remark \code{n09rel}), :441--448 (signature of the
survivor); \code{N09-RESTATE.md}:44--53.}

% =====================================================================
\subsection{Four machine-verified lemmas: two counterexamples, two certificates}
\label{sec:est-counter}

\noindent These four are model-level or arithmetic results, hence \LM: they are verified in a
solvable model or by exhaustive machine enumeration. Two of them kill implications and two certify
arithmetic, and the distinction matters, because the second pair is not evidence about Yang--Mills at
all (\S\ref{sec:est-window-fix}).

\begin{enumerate}[label=\textbf{\arabic*.},leftmargin=2.0em]
\item \textbf{What it kills:} $G''>0\Rightarrow m>0$ (``$G''$ sufficient''). \emph{Witness:} M6 has
      $G''=1/\mathrm{Var}(a)=2\omega>0$ while its spectral bottom obeys
      $m_6=\min(\omega,2\pi^2/L^2)\to0$; the finite-dimensional $2\times2$ witness is
      machine-checked, \code{sorry} $=0$, axioms
      $[\code{propext},\code{Classical.choice},\code{Quot.sound}]$. The file records that this
      witness realises the failure through a degenerate vacuum, so it kills the bare implication and
      is not claimed as an M6-class counterexample.
      \begin{sloppypar}
      \noindent\src{\code{L05Counterexamples.lean}; \code{L05-COUNTER.md};
      \code{LOOP-RUNBOOK.md}:8691--8697, :8841; \code{N09-RESTATE.md}:66 (the $m_6$ formula).}
      \end{sloppypar}
\item \textbf{What it kills:} $m>0\Rightarrow f_1^2>0$ (``$f_1^2$ necessary''). \emph{Witness:} in
      the $2$D Yang--Mills type model $[H,O]=0$ with a non-degenerate spectrum forces the
      off-diagonal element to vanish, $f_1^2=0$, while the gap stays positive; the explicit
      $2\times2$ witness is $H=\mathrm{diag}(0,1)$, $O=\mathrm{diag}(3,-1)$ with gap $1$.
      \src{\code{M03-COUNTER.lean} (\code{f1sq_eq_zero}); \code{LOOP-RUNBOOK.md}:9683--9731.}
\item \textbf{What it certifies:} the window's arithmetic at the elementary-bound step,
      $\exp(5/96)\le1.054$ and $4\exp(5/96)<7.0633$; the partial sum is $1.0534632$, the remainder
      $3.83\cdot10^{-7}$, the margin $4.6\cdot10^{-4}$.
      \src{\code{LOOP-RUNBOOK.md}:8961--8962; \code{L06-INTERVAL.md}:68--69;
      \code{L02-CBARKATO.md}:84--85.}
\item \textbf{What it certifies:} the window arithmetic as a whole --- twenty-five theorems,
      \code{sorry} $=0$, twelve axiom checks clean, including $8>7.0632853$ and the rational
      certificate. \src{\code{M02-CBAR.lean}; \code{LOOP-RUNBOOK.md}:9512--9528.}
\end{enumerate}

\noindent\textbf{How the four are used, and how they are not.} Rows 1 and 2 kill two implications
that the project had used as inference rules --- ``$G''$ is the key'', ``$f_1^2$ is the key'' --- and
they kill them on models where the implication is checkable by kernel, not by sampling. Rows 3 and 4
certify arithmetic, and arithmetic is not a statement about Yang--Mills: they certify that the numbers
in \S\ref{sec:est-window-fix} are what they are claimed to be. In particular row 4 does not make the
\emph{window} a theorem, because a numerical maximum of an explicit function is still a numerical
maximum; \LM{} is exactly the right grade for it, and it is the grade printed.

% =====================================================================
\subsection{One new \texorpdfstring{\LT}{[T]} this round: the block-Loewner order}
\label{sec:est-loewner}

\begin{claimbox}{cT}
\claimhead{Theorem}{(block-Loewner monotonicity in $\beta$).} Split the transfer kernel as
$T_\beta=D_\beta^{1/2}K_\beta D_\beta^{1/2}$ with $D_\beta=e^{-\beta S_s/2}$ and $K_\beta$ the
half-plaquette kernel. Then for $\beta>\beta'$
\begin{equation}
K_\beta\ \succeq\ K_{\beta'} ,
\label{eq:est-loewner}
\end{equation}
in the Loewner order on $L^2(dU)$: the block coefficients satisfy $c_r(\beta)\ge0$ and are
non-decreasing in $\beta$ for every representation $r$, with $c_r(\beta)=\langle\chi_r\rangle/d_r$, and
the tensor-product multiplicities entering the composition of blocks are non-negative.
\LT.
\end{claimbox}

\begin{proofx}
The order \eqref{eq:est-loewner} reduces, term by term, to $c_r(\beta)\ge c_r(\beta')$ for all $r$:
positive definiteness of $\chi_r(UV^\dagger)$ makes the block comparison a comparison of coefficients
rather than of kernels. Positivity $c_r\ge0$ and monotonicity in $\beta$ follow from the
Griffiths/Bochner route applied to $\langle\chi_r\rangle$ \cite{griffiths,ginibre,ksherman}, and the
composition of two blocks is a sum over $r_1\otimes r_2$ with non-negative multiplicities, so no
negative weight can enter the comparison.
\end{proofx}

\noindent\textbf{What it repairs, precisely.} \code{D05-FKG.md} adjudicated four readings of ``the
order of $T_\beta$'' and left exactly one slot open: at \emph{block} level the coefficient
monotonicity was known for $\mathrm{SU}(2)$ and marked \LN{} for $\mathrm{SU}(N\ge3)$
(\code{D05-FKG.md}:69, :100). \eqref{eq:est-loewner} discharges that slot, so the block-level row of
that adjudication changes from \LN{} to \LT. \src{\code{position-current.html}:380 (the round-178
record of the new \LT, item 6); \code{D05-FKG.md}:62--71.}

\smallskip
\noindent\textbf{What it does not repair.} The \emph{full-kernel} order $T_\beta\succeq T_{\beta'}$
is not established, and \eqref{eq:est-loewner} does not touch it: the factor $D_\beta$ decreases
pointwise in $\beta$, and that obstacle sits in the other slot of the splitting. And even if the
block order is granted, what it transports is the \emph{unnormalized} eigenvalue sequence --- not the
normalized kernel, not the eigenvectors, and not $F$ or $m$. The repair is a repair of one slot, not
a route to the gap. \src{\code{D05-FKG.md}:71 (the $D$-factor obstacle), :76--82 (``output zero for
$F$, for $m$, for $\mathbb G$''), :100 (the row table).}

\begin{warnbox}
\textbf{A bug, caught in the same round, that voids earlier numerical sign determinations.} A
QR-based sampler fabricates ``sign changes of $c_r$'': the sampled coefficients flip sign as a
numerical artefact of the factorization, not as a property of the kernel. Consequences, both
recorded and neither negotiable: (a) four places in the corpus state the coefficient convention in a
form that the artefact made look consistent, and those statements have to be changed; (b) several
earlier numerical conclusions of the form ``$c_r$ changes sign here'' are void and are withdrawn
pending re-derivation. The analytic statement \eqref{eq:est-loewner} is independent of the sampler,
which is why it keeps \LT{} while the numerical readings do not.
\src{\code{position-current.html}:417 (the bug entry), :399 (``two solid items: block-Loewner \LT{}
and the QR-sampler bug'').}
\end{warnbox}

% =====================================================================
\subsubsection*{What this section does not claim}

\begin{enumerate}[label=(\roman*),leftmargin=2.2em]
\item \textbf{Not Clay.} No continuum mass gap is claimed, $a\to0$ is not taken, and no statement
      here crosses to $\beta\approx2.5$. The residue and the isolation are strong-coupling
      statements at fixed spacing, and their window \eqref{eq:est-window2} is a window in the
      coupling, not in the spacing.
\item \textbf{The positivity of the mass inside the window is not ours.} For $\beta$ below the
      strong-coupling radius, $m>0$ is Osterwalder--Seiler's cluster expansion \cite{osterwalder-seiler}
      and the edge $1/48$ of the window is where Shen--Zhu--Zhu's log-Sobolev and Poincar\'e
      inequalities are proved \cite{szz}; the same conclusion in
      the polymer language is \cite{seiler,bgdf}, and for the lower-dimensional models
      \cite{gjs,fo}. What \S\ref{sec:est-i1}--\S\ref{sec:est-i2}
      add is the residue $I1$ and the isolation $I2$, i.e.\ the removal of the premise that the
      $0^{++}$ channel is strictly lightest --- and nothing more.
      \src{\code{docs/paper-v2/ch20-unconditional-i1.tex} \S D.2.}
\item \textbf{The step from \eqref{eq:est-i2} to $m_1>0$ is vacuous and is not used.} The side
      condition of the Kato-type reading is verbatim $m_1>0$, so the implication is written in
      through its own hypothesis; it is recorded as a qualification, never as a gain.
      \src{\code{docs/paper-v2/ch20-unconditional-i1.tex} \S D.1; \code{G07-KATOADJ.md} \S0 F.}
\item \textbf{Proposition C is a finite-dimensional shadow.} It covers a finite support with a fixed
      point table; it does not discharge any statement uniform in $(L,a)$ and must not be read as
      closing a gap in $\beta$. \src{\code{N09-RESTATE.md}:66--70 (\S5 item 4).}
\item \textbf{The block-Loewner order produces no lower bound.} By construction it transports
      eigenvalues, and the project's target is a lower bound on a spectral bottom; the two are
      different objects, and \S\ref{sec:est-loewner} says so in its own scope sentence.
\end{enumerate}

\noindent Everything in this section is therefore of one of two kinds: a statement whose proof is
complete at the grade printed, or a statement whose missing step is named at the place where it is
missing. There is no third kind in the register, and no claim above is graded above its weakest link.

\endgroup
